\section{Consequences and further analytic questions}
\label{end:section}

The localized energy theorem identifies an exact analytic representative
of the full zero count. Its validity through multiple zeros is a strength
for transport: the count is not lost at a collision. The same property
means that it cannot by itself exclude a collision. The regularized
measure records the divisor in a weak limit, with an explicit field-error
and localization cost. These estimates can be used on any domain for
which the boundary amplitude and the required local units are controlled.

For real heat evolution, the complete cap difference counts births only
after wall crossings and endpoint conventions are included. For complex
time, the discriminant gives the appropriate collision residue. The
cover's topology, the energy of $2\,dw$, the reciprocal differential
$dz/w$, and the physical jet phase each encode a different part of this
accounting. The exact examples show why finiteness or regularity of one
cannot be substituted for a vanishing statement about another.

The pair matrix supplies a quantitative geometric object with a precise
heat law. Its theta clock limit is analytically controlled, including
the full source tail, uniformly through $K=o(L\sqrt\kappa)$ in the
normalization of Corollary~\ref{pm:theta-growing-clock}. An absolute
physical-matrix error tending to zero follows in the smaller window
$K=o(\sqrt\kappa/L)$. These spatial estimates do not give a
growing-population dynamical law or the statistical transfer conditions
of Section~\ref{sc:statistics}. A dynamical improvement would
require an estimate for the complete exterior forcing, rather than an
unforced repulsion model. A cap estimate would additionally require the
root-derivative phase and the complete closing faces of
Section~\ref{sc:caps}.

Theorem~\ref{reg:direct-divisor} settles the fixed-disk $C^2$
weighted-divisor question with explicit constants and a sharp exponent.
Its complete cluster and boundary data remain essential.
Concrete further analytic questions are whether a coupled
source estimate controls the centered exterior residual in a useful
growing regime; whether comparable divisor estimates persist under a specified weaker
information contract or uniformly through controlled merging clusters; and
whether a full physical boundary admits a signed estimate stronger than
the identities alone. The present results settle none of those questions
by a numerical experiment or by a change of terminology. They provide
explicit formulas, errors and examples against which such estimates can
be tested.
